% ==========================================================
% Acronimi e abbreviazioni
% ==========================================================
\newacronym{api}{API}{Application Programming Interface}
\newacronym{atp}{ATP}{Available-To-Promise}
\newacronym{bom}{BOM}{Bill Of Materials}
\newacronym{ci}{CI}{Continuous Integration}
\newacronym{cli}{CLI}{Command Line Interface}
\newacronym{cpsat}{CP-SAT}{Constraint Programming -- Satisfiability}
\newacronym{erp}{ERP}{Enterprise Resource Planning}
\newacronym{grpc}{gRPC}{gRPC Remote Procedure Call}
\newacronym{gui}{GUI}{Graphical User Interface}
\newacronym{http}{HTTP}{HyperText Transfer Protocol}
\newacronym{json}{JSON}{JavaScript Object Notation}
\newacronym{kpi}{KPI}{Key Performance Indicator}
\newacronym{mrp}{MRP}{Material Requirements Planning}
\newacronym{orm}{ORM}{Object-Relational Mapping}
\newacronym{rest}{REST}{REpresentational State Transfer}
\newacronym{sql}{SQL}{Structured Query Language}
\newacronym{uml}{UML}{Unified Modeling Language}

% ==========================================================
% Glossario
% ==========================================================
\newglossaryentry{apig}{
    name={API},
    text={API},
    sort=api,
    description={Un'API (Application Programming Interface) è un insieme di procedure messe a disposizione del programmatore, tipicamente raggruppate a formare un insieme di strumenti specifici per un determinato compito. Il suo scopo è fornire un'astrazione, di norma tra l'hardware e il programmatore o tra software di basso e di alto livello, semplificando il lavoro di programmazione}
}

\newglossaryentry{umlg}{
    name={UML},
    text={UML},
    sort=uml,
    description={Nell'ingegneria del software, l'Unified Modeling Language (UML) è un linguaggio di modellazione e specifica basato sul paradigma a oggetti. UML è impiegato come ``lingua franca'' nella comunità della progettazione e della programmazione orientata agli oggetti per descrivere in modo conciso soluzioni analitiche e di progetto}
}

\newglossaryentry{solver}{
    name={solver},
    sort=solver,
    description={Componente software che risolve un problema di ottimizzazione o di soddisfacimento di vincoli, individuando una soluzione ottimale o ammissibile a partire da un modello matematico del problema}
}

\newglossaryentry{makespan}{
    name={makespan},
    sort=makespan,
    description={In ambito di scheduling, è l'istante di completamento dell'ultimo task pianificato, ovvero la durata complessiva del piano di produzione. Minimizzare il makespan significa ridurre il tempo totale necessario a completare tutte le lavorazioni}
}

\newglossaryentry{lateness}{
    name={lateness},
    sort=lateness,
    description={In ambito di scheduling, è la differenza tra l'istante di completamento di un task e la relativa data di consegna. Un valore positivo indica un ritardo rispetto alla scadenza, un valore negativo un anticipo}
}

\newglossaryentry{gantt}{
    name={diagramma di Gantt},
    text={diagramma di Gantt},
    sort=gantt,
    description={Strumento di rappresentazione grafica di un piano di attività, in cui ciascun task è disegnato come una barra orizzontale la cui posizione e lunghezza indicano rispettivamente l'istante di inizio e la durata sulla linea temporale}
}

\newglossaryentry{changeover}{
    name={changeover},
    sort=changeover,
    description={Tempo di attrezzaggio (setup) necessario a riconfigurare una macchina per passare dalla lavorazione di un articolo a quella di un articolo diverso. Costituisce un costo da minimizzare nella sequenziazione delle lavorazioni}
}

\newglossaryentry{warmstart}{
    name={warm start},
    sort=warm start,
    description={Tecnica con cui si fornisce al solver una soluzione iniziale di partenza (\textit{hint}), così da guidarne la ricerca verso le regioni più promettenti dello spazio delle soluzioni e accelerarne la convergenza}
}

\newglossaryentry{incumbent}{
    name={incumbent},
    sort=incumbent,
    description={Migliore soluzione ammissibile trovata dal solver in un dato istante del processo di ottimizzazione, soggetta a essere sostituita da soluzioni migliori nel prosieguo della ricerca}
}

\newglossaryentry{scrum}{
    name={Scrum},
    sort=scrum,
    description={Framework di sviluppo Agile che organizza il lavoro in cicli iterativi di durata fissa (\textit{sprint}), con ruoli, eventi e artefatti definiti volti a favorire la collaborazione e il rilascio incrementale del prodotto}
}

\newglossaryentry{sprint}{
    name={sprint},
    sort=sprint,
    description={Nel framework Scrum, intervallo di tempo di durata prefissata al termine del quale viene rilasciato un incremento del prodotto potenzialmente utilizzabile}
}

\newglossaryentry{kanban}{
    name={Kanban},
    sort=kanban,
    description={Metodo di gestione visuale del lavoro basato su una lavagna suddivisa in colonne che rappresentano gli stati di avanzamento delle attività, utile a limitare il lavoro in corso e a evidenziare i colli di bottiglia}
}

\newglossaryentry{gitflow}{
    name={GitFlow},
    sort=gitflow,
    description={Modello di organizzazione dei rami di un repository Git che separa lo sviluppo delle funzionalità, la preparazione delle release e le correzioni urgenti in rami dedicati (\textit{main}, \textit{develop}, \textit{feature}, \textit{release}, \textit{hotfix})}
}

\newglossaryentry{distintabase}{
    name={distinta base},
    sort=distinta base,
    description={Elenco strutturato e gerarchico dei componenti, dei semilavorati e delle materie prime necessari a realizzare un prodotto finito, con le relative quantità (in inglese \textit{Bill of Materials})}
}
